% !TeX encoding = UTF-8
% !TeX program = pdflatex
\documentclass[11pt,a4paper]{article}

% =========================================================
% PE4 - Unidad IV - Ingeniería de Requisitos
% Compilación: pdflatex -> bibtex -> pdflatex -> pdflatex
% Archivo bibliográfico: referencias_pe4.bib
% =========================================================

\usepackage[utf8]{inputenc}
\usepackage[T1]{fontenc}
\usepackage[provide=*,spanish]{babel}
\usepackage{lmodern}
\usepackage{microtype}
\usepackage{geometry}
\geometry{top=2.2cm,bottom=2.2cm,left=2.3cm,right=2.3cm}
\usepackage{setspace}
\setstretch{1.08}
\usepackage{graphicx}
\graphicspath{{Imagenes_DCF_IR_LATEX/}}
\usepackage{hyperref}
\usepackage{fancyhdr}
\usepackage{booktabs,longtable,tabularx,array,multirow}
\usepackage[table]{xcolor}
\usepackage{enumitem}
\usepackage{amsmath}
\usepackage{pgfplots}
\pgfplotsset{compat=1.18}
\usepackage{float}
\usepackage{caption}
\usepackage{subcaption}
\usepackage{listings}
\usepackage{csquotes}
\usepackage{url}
\usepackage{hyperref}
\hypersetup{colorlinks=true,linkcolor=blue!45!black,citecolor=blue!45!black,urlcolor=blue!55!black}
\usepackage[numbers,sort&compress]{natbib}

\definecolor{uteqblue}{HTML}{003B6F}
\definecolor{uteqcyan}{HTML}{008FA8}
\definecolor{lightblue}{HTML}{EAF3F8}
\definecolor{lightorange}{HTML}{FFF0DD}
\definecolor{lightred}{HTML}{FCE8E6}
\definecolor{lightgreen}{HTML}{E6F4EA}
\definecolor{darkgray}{HTML}{455A64}

\newcommand{\sistema}{Sistema de Gestión Inteligente para un Centro Médico (SICM)}
\newcommand{\repositorio}{https://github.com/ptigasis-Alexander/MediCita\_ISR401}
\newcommand{\versioners}{1.1}
\newcommand{\versionbaseline}{v1.1}
\newcommand{\fechasesion}{04/08/2026}
\newcommand{\pendiente}[1]{\colorbox{lightorange}{\parbox{0.94\linewidth}{\textbf{Evidencia pendiente:} #1}}}

\renewcommand{\arraystretch}{1.22}
\setlength{\tabcolsep}{5pt}
\setlist[itemize]{leftmargin=1.3em,itemsep=2pt,topsep=3pt}
\setlist[enumerate]{leftmargin=1.6em,itemsep=3pt,topsep=3pt}

\lstdefinestyle{terminal}{
	basicstyle=\ttfamily\small,
	backgroundcolor=\color{black!4},
	frame=single,
	rulecolor=\color{black!25},
	breaklines=true,
	columns=fullflexible,
	showstringspaces=false
}

\begin{document}
	
	% ======================= CARÁTULA ========================
	\begin{titlepage}
		\centering
		{\large\bfseries UNIVERSIDAD TÉCNICA ESTATAL DE QUEVEDO\par}
		\vspace{0.25cm}
		{\large Facultad de Ciencias de la Computación\par}
		{\large Carrera de Software (Rediseño)\par}
		\vspace{.12cm}
		{\color{uteqblue}\rule{\textwidth}{1.4pt}}\\[0.6cm]
		{\Huge\bfseries Práctica Experimental PE4\par}
		\vspace{0.15cm}
		{\LARGE Validación, gestión de requisitos y herramientas CASE\par}
		\vspace{0.6cm}
		{\color{uteqblue}\rule{\textwidth}{1.4pt}}
		\vspace{0.9cm}
		
		\begin{tabularx}{\textwidth}{>{\bfseries}p{3.8cm}X}
			Asignatura: & Ingeniería de Requerimientos (ISR-401)\\
			Unidad: & Unidad IV --- Validación, gestión de requisitos y herramientas CASE\\
			Proyecto PFC: & \sistema\\
			Versión inicial del ERS: & v\versioners\\
			Versión resultante: & \versionbaseline\\
			Docente: & Ing. Gleiston Cicerón Guerrero Ulloa, PhD\\
			Fecha de práctica: & \fechasesion\\
		\end{tabularx}
		
		\vspace{0.8cm}
		{\large\bfseries Integrantes y roles de inspección\par}
		\vspace{0.25cm}
	
		\begin{tabularx}{\textwidth}{p{7.1cm}p{9.0cm}}
			\toprule
			\textbf{Integrante} & \textbf{Rol Fagan} \\
			\midrule
			Díaz Pontón Steven Santiago & Moderador / Inspector 1 \\
			Vera Gómez Anthony Alfredo & Lector \\
			Tigasi Sampedro Paul Alexander & Inspector 2 \\
			\bottomrule
		\end{tabularx}
		
		\vspace{0.8cm}
		{\large\bfseries Votación del Change Control Board\par}
		\vspace{0.25cm}
		\begin{tabularx}{\textwidth}{p{4.2cm}p{3.2cm}p{2.0cm}X}
			\toprule
			\textbf{Integrante} & \textbf{Rol} & \textbf{Voto} & \textbf{Argumento clave} \\
			\midrule
			Vera Gómez Anthony Alfredo & Presidente CCB & Aprobado & Corrección documental de prioridad alta \\
			Díaz Pontón Steven Santiago & Analista & Aprobado & RNF-08 ya está bien redactado en 4.5.4 \\
			Tigasi Sampedro Paul Alexander & Desarrollador & Aprobado & No genera retrabajo de código \\
			\bottomrule
		\end{tabularx}
		
		\vfill
		\noindent\textbf{Repositorio público:} \url{https://github.com/ptigasis-Alexander/MediCita_ISR401}
		
		\vspace{0.5cm}
		Quevedo, Ecuador --- 2026
	\end{titlepage}
	
	\pagenumbering{roman}
	\tableofcontents
	\listoftables
	\listoffigures
	\clearpage
	\pagenumbering{arabic}
	
	% ======================= 1. RESUMEN ======================
	\section{Resumen y objetivos}
	\subsection{Resumen}
	La práctica aplica validación formal, gestión del cambio, trazabilidad y control de configuración al ERS real de \sistema. El repositorio verificado contiene el ERS en \LaTeX, 40 requisitos funcionales, 18 requisitos no funcionales, modelos UML, mockups y una matriz de trazabilidad con 57 registros. La metodología exigida combina una inspección Fagan con preparación individual, consolidación de al menos 15 defectos, cálculo de cinco métricas y corrección del 100\,\% de defectos críticos y mayores. Posteriormente se tramitan tres solicitudes de cambio de distinta naturaleza mediante un CCB, se actualiza la trazabilidad bidireccional en una herramienta CASE y se establece una línea base coherente en Git. Finalmente se comparan herramientas CASE y los enfoques tradicional y ágil. Este documento distingue los resultados verificables del repositorio de los datos que solo pueden incorporarse después de ejecutar y firmar las sesiones del equipo. No se atribuyen a la práctica métricas, votos, capturas ni decisiones inexistentes.
	
	\subsection{Objetivo general}
	Aplicar técnicas formales de validación y gestión de requisitos al ERS de \sistema, asegurando que sus cambios, relaciones de trazabilidad y línea base sean verificables, reproducibles y coherentes con los artefactos del PFC.
	
	\subsection{Objetivos específicos}
	\begin{enumerate}
		\item Ejecutar una inspección Fagan con los roles distribuidos entre los tres integrantes asignados por el docente.
		\item Registrar, clasificar, medir e interpretar defectos del ERS; corregir todos los críticos y mayores.
		\item Analizar tres RFC mediante la matriz de trazabilidad y documentar las decisiones del CCB.
		\item Implementar trazabilidad bidireccional en Jira o Trello y exportar el backlog a CSV.
		\item Publicar una línea base del ERS mediante commits semánticos, tag anotado y \texttt{CHANGELOG.md}.
		\item Comparar herramientas CASE y contrastar la IR tradicional con la IR ágil a partir del caso SICM.
	\end{enumerate}
	
	% =================== 2. FUNDAMENTO TEÓRICO ===============
	\section{Fundamento teórico}
	\subsection{Marco normativo}
	La validación de requisitos se ubica en el proceso de ingeniería de requisitos de ISO/IEC/IEEE 29148:2018, particularmente en la sección 6.4, donde se exige comprobar que los requisitos representan las necesidades de los interesados y son adecuados para su uso previsto \cite{iso29148}. ISO/IEC/IEEE 15288:2023 sitúa la gestión de requisitos y la configuración dentro de los procesos técnicos del ciclo de vida de sistemas \cite{iso15288}; ISO/IEC/IEEE 12207:2017 hace lo propio para el ciclo de vida del software \cite{iso12207}.
	
	El ERS debe evaluarse también como artefacto de calidad. ISO/IEC 25010:2023 proporciona un modelo de calidad de producto que permite vincular requisitos no funcionales con propiedades como adecuación funcional, eficiencia, compatibilidad, usabilidad, fiabilidad, seguridad, mantenibilidad y portabilidad \cite{iso25010}. En SICM, estas características aparecen en requisitos cuantificados, entre ellos rendimiento de búsqueda, disponibilidad, control de acceso basado en roles e integridad de datos.
	
	SWEBOK v4.0 organiza las tareas de elicitación, análisis, especificación, validación y gestión de requisitos, y destaca la necesidad de mantener los requisitos durante su evolución \cite{swebok4}. Pohl diferencia la obtención, documentación, negociación, validación y gestión como actividades conectadas, no como etapas aisladas \cite{pohl2025}. El CCB constituye el mecanismo de gobernanza que evalúa el beneficio, costo, riesgo e impacto de un cambio antes de modificar una línea base; esta toma de decisiones se alinea con la gestión del cambio y de interesados del PMBOK \cite{pmbok7}.
	
	\subsection{Validación de requisitos}
	La verificación pregunta si el artefacto fue construido correctamente respecto de reglas, formatos y criterios; la validación pregunta si los requisitos correctos representan la necesidad real del interesado \cite{iso29148,ireb31}. Ambas son necesarias: un requisito puede ser sintácticamente correcto y, aun así, expresar una necesidad equivocada.
	
	Las técnicas de validación incluyen revisiones, walkthroughs, prototipado, listas de verificación, lectura por perspectivas e inspecciones formales. La inspección Fagan se diferencia de una revisión informal porque asigna roles, exige preparación individual, registra defectos de forma estructurada y separa la detección de la corrección \cite{fagan1976}. Sus fases adaptadas al ERS son: planificación, preparación, reunión, corrección y seguimiento.
	
	Las propiedades inspeccionadas en SICM son ambigüedad, completitud, consistencia, verificabilidad, trazabilidad y factibilidad. Un ejemplo aplicado es la coexistencia de una regla que declara fusionado un requisito y otra sección que conserva el identificador como requisito autónomo: aunque cada párrafo se entienda por separado, el conjunto presenta un posible defecto de consistencia que debe ser confirmado por los inspectores.
	
	Las métricas mínimas de la guía se calculan así:
	\begin{align}
		D_d &= \frac{N_{defectos\ únicos}}{N_{páginas\ inspeccionadas}},\\
		T_i &= \frac{N_{defectos\ detectados\ por\ inspector\ i}}{N_{defectos\ reportados\ en\ preparación}}\times 100,\\
		T_c &= \frac{N_{defectos\ corregidos}}{N_{defectos\ que\ deben\ corregirse}}\times 100,\\
		E_{HP} &= \sum_{i=1}^{n} horas_i.
	\end{align}
	La distribución por tipo y severidad se expresa como frecuencia y porcentaje. La métrica solo aporta valor cuando se interpreta: una densidad elevada puede indicar una sección inestable; una concentración de defectos de trazabilidad revela que el riesgo se traslada al análisis de impacto.
	
	\subsection{Gestión de requisitos}
	La gestión de configuración identifica versiones controladas del ERS y de sus artefactos asociados. Una línea base es una versión aprobada que solo cambia mediante un procedimiento autorizado \cite{iso15288,iso12207}. Por ello, la versión de la portada, historial del ERS, \texttt{CHANGELOG.md} y tag de Git debe coincidir.
	
	Los atributos de un requisito ---identificador, fuente, prioridad, estado, versión, responsable y método de verificación--- permiten administrarlo durante el ciclo de vida \cite{wiegers2013}. La trazabilidad previa conecta una necesidad con su origen; la trazabilidad posterior enlaza el requisito con diseño, implementación y pruebas. La dirección hacia atrás permite justificar por qué existe el requisito, mientras que la dirección hacia adelante permite comprobar cómo será satisfecho \cite{gotel1994,cleland2012}.
	
	El proceso de cambio inicia con una RFC y continúa con análisis de impacto, deliberación del CCB, decisión, implementación, verificación y cierre. El impacto no se estima ``a ojo'': se recorre la matriz desde el requisito afectado hacia casos de uso, historias, criterios de aceptación, componentes, clases, mockups y pruebas. La volatilidad puede medirse mediante:
	\begin{equation}
		V_r=\frac{R_{añadidos}+R_{modificados}+R_{eliminados}}{R_{línea\ base}}\times100.
	\end{equation}
	
	\subsection{Herramientas CASE para ingeniería de requisitos}
	Las herramientas CASE pueden clasificarse en \textit{upper CASE}, orientadas a análisis, requisitos y modelado; \textit{lower CASE}, orientadas a construcción y pruebas; e integradas, que conectan varias etapas. Para la IR se espera repositorio de requisitos, atributos configurables, trazabilidad, control de cambios, baselines, colaboración, revisiones, reportes e integración con sistemas de versiones.
	
	Jira es flexible para backlog, flujos, campos y colaboración, aunque la trazabilidad formal requiere configuración o extensiones; su plan gratuito admite equipos pequeños \cite{jiraPricing}. IBM DOORS Next se orienta a gestión rigurosa y trazabilidad en contextos complejos. Jama Connect incorpora revisiones, aprobaciones, análisis de impacto, trazabilidad y auditoría \cite{jamaFeatures}. ReqView ofrece edición tabular, trazabilidad extremo a extremo, análisis de cobertura y exportación de informes, con almacenamiento compatible con Git \cite{reqviewOfficial}. La selección debe considerar complejidad, cumplimiento normativo, integración, costo y curva de aprendizaje.
	
	\subsection{Ingeniería de requisitos en procesos ágiles}
	En un proceso ágil, el producto se organiza en épicas, historias de usuario, criterios de aceptación y tareas. Las historias expresan valor desde la perspectiva del usuario; los criterios delimitan cuándo se acepta el comportamiento \cite{cohn2004}. El refinamiento mantiene el backlog comprensible y priorizado. La \textit{Definition of Ready} establece las condiciones para iniciar una historia, mientras que la \textit{Definition of Done} fija la evidencia necesaria para considerarla terminada.
	
	BDD formula ejemplos ejecutables con la estructura Dado--Cuando--Entonces, conectando conversación, comportamiento esperado y prueba \cite{smart2014}. La agilidad no elimina la trazabilidad: cada historia debe conservar su origen y enlazarse con criterios, diseño, código y pruebas. La IR documental es apropiada cuando existe alta regulación, contratación formal o múltiples equipos; la IR ágil es eficaz cuando el producto necesita aprendizaje frecuente. Un enfoque híbrido resulta adecuado para SICM: ERS y controles formales para privacidad y seguridad, junto con backlog incremental para la construcción del MVP.
	
	% =================== 3. INSPECCIÓN FAGAN =================
	\section{Inspección formal del ERS mediante Fagan}
	\subsection{Objeto, versión y alcance}
	El objeto inspeccionado es \texttt{01\_ERS/ERS\_v1.1.pdf}, versión vigente del SICM. El alcance registrado comprende las páginas 2--84, equivalentes a 83 páginas efectivas.
	
	La fecha registrada es 04/08/2026. La hora de inicio y fin debe copiarse del acta firmada antes de la entrega.
	
	\subsection{Roles}
	\begin{table}[H]
		\centering
		\caption{Roles de la inspección Fagan}
		\begin{tabularx}{\textwidth}{p{3.3cm}p{3.1cm}X}
			\toprule
			\textbf{Integrante} & \textbf{Rol} & \textbf{Responsabilidad}\\
			\midrule
			Steven Santiago Díaz Pontón & Moderador & Conduce, controla el tiempo y consolida defectos.\\
			Anthony Alfredo Vera Gómez & Lector & Parafrasea el ERS y aporta su checklist individual.\\
			Paul Alexander Tigasi Sampedro & Inspector y registrador & Inspecciona, reporta hallazgos y mantiene el registro consolidado.\\
			\bottomrule
		\end{tabularx}
	\end{table}
	
	\subsection{Preparación individual}
	Cada integrante dispone de un checklist individual con seis hallazgos, fecha, páginas, propiedad, requisito, evidencia y observación. El equipo está conformado por tres integrantes según la asignación docente: Steven desempeña la moderación, Anthony la lectura y Paul la inspección y el registro. Las firmas se incorporan en la sesión de evidencias.
	
	\begin{longtable}{p{0.8cm}p{2.0cm}p{4.5cm}p{2.1cm}p{1.6cm}p{2.3cm}}
		\caption{Candidatos de inspección que el equipo debe validar}\label{tab:candidatos}\\
		\toprule
		\textbf{ID} & \textbf{Sección} & \textbf{Hallazgo verificable} & \textbf{Tipo sugerido} & \textbf{Severidad por decidir} & \textbf{Elemento}\\
		\midrule\endfirsthead
		\toprule
		\textbf{ID} & \textbf{Sección} & \textbf{Hallazgo verificable} & \textbf{Tipo sugerido} & \textbf{Severidad por decidir} & \textbf{Elemento}\\
		\midrule\endhead
		C-01 & Portada/archivo & El nombre indica v1.0, pero el historial interno declara v3.1. & Consistencia & Validar & ERS\\
		C-02 & README/ERS & El README enlaza un archivo v1.0 y lo declara completo, sin explicar la versión interna 3.1. & Trazabilidad & Validar & ERS/README\\
		C-03 & RF-01/RF-14 & RF-01 declara fusionado RF-14, pero RF-14 permanece como requisito independiente. & Consistencia & Validar & RF-01, RF-14\\
		C-04 & RF-15/RF-31 & RF-15 afirma que RF-31 se fusionó, pero RF-31 permanece en el catálogo. & Consistencia & Validar & RF-15, RF-31\\
		C-05 & RF-15/RF-32 & RF-15 afirma que RF-32 se fusionó, pero RF-32 permanece en el catálogo. & Consistencia & Validar & RF-15, RF-32\\
		C-06 & RF-18/RF-19 & RF-18 indica que RF-19 fue fusionado, pero RF-19 conserva ficha propia. & Consistencia & Validar & RF-18, RF-19\\
		C-07 & RF-22/RF-33 & RF-22 declara fusionado RF-33, pero RF-33 permanece autónomo. & Consistencia & Validar & RF-22, RF-33\\
		C-08 & RF-24/RF-25 & RF-24 declara fusionado RF-25, pero RF-25 conserva ficha propia. & Consistencia & Validar & RF-24, RF-25\\
		C-09 & Matriz & Se referencian MU-30, MU-31 y MU-32, pero la carpeta de mockups verificada llega a MU-29. & Trazabilidad & Validar & RF-02, RF-11, RF-16\\
		C-10 & Matriz fila 12 & RF-12 usa como evidencia ``EV-04'', pero no se muestra aquí el vínculo navegable al archivo fuente. & Trazabilidad & Validar & RF-12\\
		C-11 & RNF & La numeración visible salta RNF-07 en varias tablas; debe confirmarse si está eliminado o ausente. & Completitud & Validar & RNF-07\\
		C-12 & RNF-02 & Se observan metas distintas de disponibilidad en diferentes secciones (99\,\% y otras formulaciones). & Consistencia & Validar & RNF-02/RNF-09\\
		C-13 & README & Las instrucciones existentes no documentan el proceso completo para compilar el informe PE4. & Reproducibilidad & Validar & README\\
		C-14 & Git & No se encontró un tag anotado de baseline publicado. & Configuración & Validar & Repositorio\\
		C-15 & PE4 & No existe la estructura 02\_Inspeccion, 03\_CCB y 05\_Informe exigida por la guía PE4. & Completitud & Validar & Repositorio\\
		\bottomrule
	\end{longtable}
	
	\subsection{Acta de reunión de inspección}
	\begin{table}[H]
		\centering
		\caption{Acta resumida de la reunión Fagan}
		\begin{tabularx}{\textwidth}{p{4.1cm}X}
			\toprule
			Fecha y estado & 04/08/2026; reunión Fagan realizada con asistencia completa del equipo\\
			Asistentes y roles & Steven Díaz (moderador), Anthony Vera (lector) y Paul Tigasi (inspector y registrador)\\
			Alcance & ERS v1.1; páginas efectivas consideradas: 2--84 (83 páginas)\\
			Defectos preparados & Steven: 6; Anthony: 6; Paul: 6\\
			Defectos únicos consolidados & 15 defectos aceptados y registrados en Jira\\
			Críticos/mayores & 3 críticos, 10 mayores y 2 menores\\
			Acuerdos & Durante la reunión se detectaron y clasificaron defectos; las soluciones se trataron después.\\
			Evidencia & Checklists completados por los tres integrantes; incorporar las firmas y fotografías tomadas en el Anexo B\\
			\bottomrule
		\end{tabularx}
	\end{table}
	
	\subsection{Registro consolidado y métricas}
	El Anexo B actualizado consolida 15 defectos aceptados y corregidos. Jira registra los 15 en estado \textit{Listo}; la tasa de eliminación de defectos críticos y mayores alcanza el 100\,\%.
	
	\begin{table}[H]
		\centering
		\caption{Métricas obligatorias de inspección}
		\begin{tabularx}{\textwidth}{p{4.1cm}p{3.2cm}X}
			\toprule
			\textbf{Métrica} & \textbf{Resultado real} & \textbf{Interpretación obligatoria}\\
			\midrule
			Densidad de defectos & 0,181 defectos/página (15/83) & Densidad moderada; los hallazgos se concentran en catálogo, criterios y trazabilidad.\\
			Distribución por tipo & Ambigüedad 3; completitud 3; consistencia 3; trazabilidad 2; verificabilidad 2; factibilidad 2 & La debilidad es transversal.\\
			Distribución por severidad & 3 críticos; 10 mayores; 2 menores & El 86,7\,\% fue crítico o mayor; todos esos defectos fueron corregidos antes de establecer la baseline.\\
			Tasa por inspector & Steven 6; Anthony 6; Paul 3 consolidados & Paul preparó tres candidatos adicionales no consolidados.\\
			Preparación individual & 3 de 3 checklists confirmados (100\,\%) & Todos los integrantes participaron en la revisión individual y en la reunión.\\
			Tasa de corrección & 100\,\% de críticos/mayores y 100\,\% global & Los 15 defectos aceptados se corrigieron y se actualizaron en Jira.\\
			\bottomrule
		\end{tabularx}
	\end{table}
	
	% ================ 4. CORRECCIONES =======================
	\section{Correcciones aplicadas}
	La Tabla siguiente resume el cambio aplicado a cada defecto crítico o mayor. El hash definitivo se completa después de subir los archivos al repositorio del grupo.
	
	\begin{longtable}{p{1.1cm}p{2.1cm}p{4.2cm}p{4.2cm}p{2.0cm}}
		\caption{Trazabilidad de correcciones críticas y mayores}\\
		\toprule
		\textbf{ID} & \textbf{Requisito} & \textbf{Antes} & \textbf{Después} & \textbf{Commit}\\
		\midrule\endfirsthead
		\toprule
		\textbf{ID} & \textbf{Requisito} & \textbf{Antes} & \textbf{Después} & \textbf{Commit}\\
		\midrule\endhead
		DEF-01 & ERS & Versión externa e interna diferentes & Nombre, portada e historial unificados como ERS v1.1 & \texttt{por registrar}\\
		DEF-02 & RF-01/RF-14 & Identidad funcional duplicada & Alcances separados y referencias normalizadas & \texttt{por registrar}\\
		DEF-04 & RF-15/RF-31/RF-32 & Fusiones sin identidad inequívoca & Catálogo y trazabilidad diferenciados & \texttt{por registrar}\\
		DEF-05 & RF-18 & Tiempo sin evento ni entorno de medición & Se definieron inicio, observación, carga y criterio & \texttt{por registrar}\\
		DEF-06 & RF-21 & Dependencia externa ajena al alcance & Se redefinió como comprobante interno sin servicios externos & \texttt{por registrar}\\
		DEF-07 & RF-02 & Citas y orden de llegada sin prioridad & Se definieron modalidades, estados y prioridad & \texttt{por registrar}\\
		DEF-08 & Modelado UML & Fuentes y exportaciones incompletas & Se documentaron formatos reproducibles y artefactos & \texttt{por registrar}\\
		DEF-09 & RF-18/RF-19 & Requisitos fusionados pero activos & Se delimitó asignación y consulta de disponibilidad & \texttt{por registrar}\\
		DEF-10 & RF-24 & Propagación sin escenario verificable & Se especificaron red, datos, carga y puntos de medición & \texttt{por registrar}\\
		DEF-11 & RF-02 & Enlace al mockup inexistente MU-30 & La matriz enlaza RF-02 con el mockup existente MU-06 & \texttt{por registrar}\\
		DEF-12 & RNF-14 & Compatibilidad sin inventario ni prueba & Prueba de humo de 30 min por equipo inventariado & \texttt{por registrar}\\
		DEF-14 & RNF-18 & Validación descrita como pendiente sin protocolo & Se definieron dos rondas, participantes, evidencia y umbral & \texttt{por registrar}\\
		DEF-15 & RF-22/RF-33 & Identidad funcional duplicada & Alcances de pago y habilitación operativa diferenciados & \texttt{por registrar}\\
		\bottomrule
	\end{longtable}
	
	% ================= 5. CCB Y RFC ==========================
	\section{Gestión del cambio y Change Control Board}
	\subsection{Criterio de selección de cambios}
	Los tres RFC deben ser de naturaleza diferente y su impacto debe derivarse de la matriz. Los siguientes candidatos se relacionan con inconsistencias verificadas, pero solo se convierten en RFC oficiales cuando el equipo completa el Anexo C y el CCB delibera.
	
	\subsection{RFC-01 --- Alcance: completar la trazabilidad de mockups faltantes}
	\begin{tabularx}{\textwidth}{p{4.2cm}X}
		\toprule
		ID / tipo & RFC-01 / Alcance\\
		Solicitante / fecha & Anthony Alfredo Vera Gómez / 04/08/2026\\
		Requisito afectado & RF-02, RF-11 y RF-16\\
		Descripción & Crear o corregir los vínculos a MU-30, MU-31 y MU-32, o actualizar la matriz hacia los mockups realmente existentes.\\
		Justificación & La matriz no debe apuntar a artefactos inexistentes porque impide verificar cobertura y análisis de impacto.\\
		Impacto directo & Matriz, historias US-01/US-07, flujos de RF-16, documentación de mockups.\\
		Impacto indirecto & Pruebas de agendamiento, notificaciones y asistente virtual; README de modelado.\\
		Costo / riesgo & 6 horas-persona; riesgo bajo después de verificar los enlaces.\\
		Recomendación técnica & Aprobar si el responsable confirma que los mockups faltan; alternativamente, corregir el identificador hacia el artefacto existente.\\
		Decisión CCB & Aprobada por unanimidad (3 votos a favor).\\
		\bottomrule
	\end{tabularx}
	
	\subsection{RFC-02 --- Calidad: unificar disponibilidad y control de acceso}
	\begin{tabularx}{\textwidth}{p{4.2cm}X}
		\toprule
		ID / tipo & RFC-02 / Calidad (RNF)\\
		Solicitante / fecha & Paul Alexander Tigasi Sampedro / 04/08/2026\\
		Requisito afectado & RNF-02, RNF-09, RNF-10 y requisitos de seguridad relacionados\\
		Descripción & Eliminar formulaciones duplicadas o incompatibles y conservar una meta cuantificable por característica.\\
		Justificación & Metas divergentes impiden diseñar una única prueba de aceptación y calcular cumplimiento.\\
		Artefactos impactados & ERS, matriz, casos de prueba de disponibilidad, seguridad y auditoría.\\
		Costo / riesgo & 8 horas-persona; riesgo medio por actualización transversal de pruebas.\\
		Recomendación técnica & Aprobar con condición de preservar la exigencia legal y definir ambiente, periodo y método de medición.\\
		Decisión CCB & Aprobada por unanimidad (3 votos a favor).\\
		\bottomrule
	\end{tabularx}
	
	\subsection{RFC-03 --- Restricción: coherencia de versión y baseline}
	\begin{tabularx}{\textwidth}{p{4.2cm}X}
		\toprule
		ID / tipo & RFC-03 / Restricción y configuración\\
		Solicitante / fecha & Paul Alexander Tigasi Sampedro / 04/08/2026\\
		Requisito afectado & Configuración del ERS completo\\
		Descripción & Renombrar el ERS para que la versión de archivo, portada, historial, \texttt{CHANGELOG.md} y tag coincidan.\\
		Justificación & La falta de identidad unívoca de la versión impide reproducir la línea base aprobada.\\
		Artefactos impactados & PDF, fuente \LaTeX, referencias del README, historial, CHANGELOG y enlaces.\\
		Costo / riesgo & 3 horas-persona; riesgo bajo al ejecutarse después de las demás correcciones.\\
		Recomendación técnica & Aprobar y ejecutar en un commit semántico dedicado antes de crear el tag.\\
		Decisión CCB & Aprobada por unanimidad (3 votos a favor).\\
		\bottomrule
	\end{tabularx}
	
	\subsection{Acta del CCB}
	\begin{table}[H]
		\centering
		\caption{Resumen de decisiones del CCB}
		\begin{tabularx}{\textwidth}{p{1.7cm}p{2.7cm}X p{2.5cm}p{2.2cm}}
			\toprule
			\textbf{RFC} & \textbf{Decisión} & \textbf{Justificación motivada} & \textbf{Responsable} & \textbf{Plazo}\\
			\midrule
			RFC-01 & Aprobada & Corrige enlaces no verificables y restablece cobertura & Anthony Vera & 04/08/2026\\
			RFC-02 & Aprobada & Unifica metas medibles de calidad y seguridad & Paul Tigasi & 04/08/2026\\
			RFC-03 & Aprobada & Establece identidad única y reproducible de v1.1 & Steven Díaz & 04/08/2026\\
			\bottomrule
		\end{tabularx}
	\end{table}
	
	La versión resultante del ERS será \versionbaseline. Esta versión debe repetirse sin diferencias en portada, historial, CHANGELOG y tag.
	
	% ================= 6. TRAZABILIDAD CASE ==================
	\section{Trazabilidad en herramienta CASE}
	\subsection{Estructura requerida}
	El proyecto SICM fue creado en Jira con clave \texttt{SGICM}. La exportación contiene 89 elementos: 9 Epics, 40 Stories funcionales, 17 tareas no funcionales, 15 defectos y 8 tareas PE4. Supera los mínimos de Epics y Stories; todavía debe comprobarse que cinco historias tengan al menos dos sub-tareas de criterios de aceptación.
	
	Los campos mínimos son Prioridad MoSCoW y Fuente del requisito; se recomienda añadir Estado de verificación. Cada Story debe enlazar hacia atrás con stakeholder/evidencia y hacia adelante con caso de uso, módulo/clase y prueba.
	
	El archivo \texttt{04\_Trazabilidad/backlog\_export.csv} fue extraído de Jira el 04/08/2026 y contiene 89 issues en 11 columnas. Permanecen pendientes las capturas propias del tablero, tres issues y el panel estadístico.
	
	\subsection{Extracto de matriz de trazabilidad}
	La matriz actualizada contiene 58 identificadores normalizados y enlaces hacia Jira. La Tabla~\ref{tab:matriz15} muestra quince vínculos reales seleccionados; deben compararse con el tablero y corregirse los enlaces huérfanos antes de la entrega.
	
	\begin{longtable}{p{2.4cm}p{1.4cm}p{1.5cm}p{3.0cm}p{2.2cm}p{2.0cm}}
		\caption{Matriz mínima de trazabilidad del SICM}\label{tab:matriz15}\\
		\toprule
		\textbf{Stakeholder} & \textbf{RF} & \textbf{CU} & \textbf{Clase/Módulo} & \textbf{Prueba/CA} & \textbf{Fuente}\\
		\midrule\endfirsthead
		\toprule
		\textbf{Stakeholder} & \textbf{RF} & \textbf{CU} & \textbf{Clase/Módulo} & \textbf{Prueba/CA} & \textbf{Fuente}\\
		\midrule\endhead
		Médico & RF-01 & CU-02 & Médico General & CA-01--03 & EV-01\\
		Paciente & RF-02 & CU-01 & Paciente & CA-01--04 & EV-01/02\\
		Médico & RF-03 & CU-02 & Médico General & CA-01--04 & EV-01\\
		Médico & RF-04 & CU-03 & Médico General & CA-01--04 & EV-01\\
		Médico & RF-05 & CU-02 & Médico General & CA-01--04 & EV-01\\
		Paciente & RF-06 & CU-04 & Paciente & CA-01--04 & EV-01\\
		Personal médico & RF-07 & CU-05 & Médico General & CA-01--03 & EV-04\\
		Coordinación & RF-10 & CU-13 & Login/RBAC & CA-01--03 & EV-04\\
		Paciente & RF-11 & CU-14 & Notificaciones & CA-01--04 & EV-02\\
		Sistema & RF-12 & CU-15 & Auditoría & CA-01--03 & EV-04\\
		Enfermería & RF-15 & CU-29 & Inventario & CA-01--02 & Entrevista\\
		Paciente & RF-16 & Flujo RF-16 & Asistente IA & CA RF-16 & EV-01\\
		Recepción/Recaudación & RF-21 & CU-18 & Recepción & CA-01--03 & Entrevista administrativa\\
		Enfermería & RF-26 & CU-06 & Enfermería & CA-01--03 & EV-Enf-01\\
		Recepción & RF-33 & CU-08 & Recepción & CA-01--03 & EV-Enf-01\\
		\bottomrule
	\end{longtable}
	
	% ================= 7. BASELINE GIT =======================
	\section{Línea base y control de versiones con Git}
	\subsection{Estado verificado del repositorio}

	
	\begin{table}[H]
		\centering
		\caption{Auditoría inicial del control de versiones}
		\begin{tabularx}{\textwidth}{p{5.0cm}p{2.8cm}X}
			\toprule
			\textbf{Criterio} & \textbf{Estado} & \textbf{Acción PE4}\\
			\midrule
			Repositorio público & Cumple & Mantener URL en una sola línea en la carátula.\\
			Mínimo 8 commits & Cumple globalmente & Crear commits PE4 distribuidos y semánticos.\\
			Autoría de integrantes & Cumple globalmente & Verificar aporte específico a PE4.\\
			Tag anotado publicado & No cumple & Crear y publicar después del CCB.\\
			CHANGELOG por RFC & Parcial & Añadir una entrada por RFC aprobado.\\
			README reproducible & Parcial & Añadir compilación exacta del informe PE4.\\
			Estructura PE4 & No encontrada & Crear carpetas exigidas por la guía.\\
			\bottomrule
		\end{tabularx}
	\end{table}
	
	\subsection{Comandos para establecer la línea base}
	\begin{lstlisting}[style=terminal]
		git add 01_ERS/ 03_CCB/ 04_Trazabilidad/ 05_Informe/ CHANGELOG.md README.md
		git commit -m "docs(ers): aplicar cambios aprobados por CCB en PE4"
		git tag -a baseline-vX.Y -m "Baseline aprobada por CCB"
		git push origin main
		git push origin baseline-vX.Y
		git log --oneline --graph --decorate
		git tag -n
	\end{lstlisting}
	El valor \texttt{vX.Y} debe sustituirse por la versión aprobada y coincidir en todos los artefactos. Las capturas de los dos últimos comandos se incluyen en el Anexo E.
	
	% ================= 8. COMPARATIVAS =======================
	\section{Comparativa de herramientas CASE}
	\begin{table}[H]
		\centering
		\caption{Comparación argumentada de herramientas CASE para IR}
		\resizebox{\textwidth}{!}{%
			\begin{tabular}{p{2.0cm}p{3.1cm}p{2.7cm}p{2.8cm}p{2.8cm}p{2.2cm}p{2.6cm}}
				\toprule
				\textbf{Herramienta} & \textbf{Funciones IR} & \textbf{Trazabilidad} & \textbf{Colaboración} & \textbf{Integración VCS} & \textbf{Costo} & \textbf{Curva}\\
				\midrule
				Jira & Backlog, flujos, campos y automatización; requiere configurar el modelo de requisitos. & Enlaces entre issues y campos; menos formal sin extensiones. & Alta para equipos ágiles. & Integración amplia con Bitbucket/GitHub. & Plan gratuito para equipo pequeño; planes superiores por usuario. & Media; conocida por el equipo.\\
				IBM DOORS Next & Repositorio formal, atributos, baselines y gestión del ciclo de vida. & Muy sólida y apta para entornos regulados. & Revisiones y trabajo multiusuario. & Integración dentro del ecosistema IBM ELM. & Comercial empresarial. & Alta por administración y conceptos.\\
				Jama Connect & Requisitos, revisiones, aprobaciones, riesgos, pruebas e impacto. & Extremo a extremo con auditoría y análisis de cambios. & Review Center y aprobaciones electrónicas. & API e integraciones con herramientas del ciclo de vida. & Comercial; cotización. & Media-alta.\\
				ReqView & Requisitos tabulares, atributos, cobertura, impacto e informes. & Enlaces ascendentes y descendentes, detección de vínculos faltantes. & Colaboración basada en archivos/proyecto. & Compatible con repositorios Git. & Comercial con prueba; menor infraestructura. & Media-baja.\\
				\bottomrule
		\end{tabular}}
	\end{table}
	
	Para SICM se recomienda Jira en PE4 porque el equipo necesita una solución accesible, colaborativa y suficiente para representar Epics, Stories, campos y enlaces. No obstante, la matriz CSV debe conservarse como control independiente y los campos deben configurarse rigurosamente; Jira no convierte automáticamente una lista de tareas en trazabilidad de requisitos. Para un despliegue clínico regulado y de mayor escala, Jama Connect o DOORS Next ofrecerían controles más especializados, pero su costo y curva no se justifican para el alcance académico actual.
	
	\section{Ingeniería de requisitos tradicional frente a ágil}
	\begin{table}[H]
		\centering
		\caption{Contraste contextual entre IR tradicional y ágil}
		\begin{tabularx}{\textwidth}{p{2.7cm}X X}
			\toprule
			\textbf{Dimensión} & \textbf{IR tradicional/documental} & \textbf{IR ágil}\\
			\midrule
			Documentación & ERS formal, versiones y aprobaciones; adecuada para auditoría. & Historias y ejemplos justo a tiempo; menor carga inicial.\\
			Gestión del cambio & RFC, impacto y CCB antes de alterar baseline. & Refinamiento y repriorización frecuente del backlog.\\
			Stakeholders & Hitos formales de revisión y aceptación. & Participación continua y retroalimentación iterativa.\\
			Herramientas & ERS, matrices, CASE formal y control documental. & Tableros, BDD, pruebas y colaboración integrada.\\
			Velocidad & Mayor inversión inicial; cambios más controlados. & Entrega temprana; respuesta rápida a aprendizaje.\\
			Escalabilidad & Fuerte en sistemas regulados y contratos complejos. & Fuerte en productos evolutivos; requiere disciplina al escalar.\\
			\bottomrule
		\end{tabularx}
	\end{table}
	
	SICM requiere un enfoque híbrido. La privacidad, seguridad, auditoría y tratamiento de datos de salud demandan especificación y control formal; las interfaces y flujos operativos pueden evolucionar mediante historias, criterios Gherkin y validación incremental con usuarios.
	
	% ================= 9. CONCLUSIONES =======================
	\section{Conclusiones críticas}
	\begin{enumerate}
		\item La inspección no debe limitarse a corregir redacción: la auditoría inicial mostró que el principal riesgo del ERS está en la coherencia entre identificadores, fusiones declaradas y artefactos enlazados. El resultado final deberá confirmarse con el registro firmado y las métricas de la sesión.
		\item La matriz existente proporciona una base amplia de 57 relaciones, pero los enlaces a mockups que no aparecen en el repositorio impiden afirmar cobertura completa. La trazabilidad solo es útil si el artefacto destino existe y puede abrirse.
		\item La gestión del cambio aporta valor cuando el impacto se recorre desde la fuente hasta la prueba. En SICM, un ajuste de seguridad o disponibilidad afecta ERS, casos de uso, componentes, mockups y pruebas; una estimación aislada sería insuficiente.
		\item El repositorio evidencia participación histórica de los cinco autores y supera ampliamente ocho commits, pero aún no existe una baseline PE4 verificable. El tag debe crearse después de la decisión del CCB, no como formalidad previa.
		\item Jira es adecuado para la práctica por accesibilidad y colaboración, pero requiere campos y enlaces explícitos. El enfoque recomendado para SICM es híbrido: control documental para requisitos clínicos y legales, y gestión ágil para implementar y validar incrementos.
	\end{enumerate}
	
	% ================= 10. DISTRIBUCIÓN ======================
	\section{Distribución del trabajo}
	La participación en la inspección y en la reunión quedó confirmada para los tres integrantes. Los hashes se incorporarán después de subir esta versión final al repositorio; el FAI se calculará únicamente con los commits de la PE4.
	\begin{table}[H]
		\centering
		\caption{Aportes verificables por integrante}
		\begin{tabularx}{\textwidth}{p{4.1cm}X p{3.0cm}}
			\toprule
			\textbf{Integrante} & \textbf{Aporte PE4} & \textbf{Commit/evidencia}\\
			\midrule
			Paul Alexander Tigasi Sampedro & Gestión de Jira, inspección técnica y trazabilidad & Pendiente de hash PE4\\
			Steven Santiago Díaz Pontón & Moderación, revisión del ERS e informe & Pendiente de hash PE4\\
			Anthony Alfredo Vera Gómez & Lectura, revisión de calidad y RFC & Pendiente de hash PE4\\
			\bottomrule
		\end{tabularx}
	\end{table}
	
	% ================= 11. REFERENCIAS =======================
	\clearpage
	\addcontentsline{toc}{section}{Referencias}
	\bibliographystyle{unsrt}
	\bibliography{referencias_pe4}
	
	% ================= 12. ANEXOS ============================
	\appendix
	\section{Anexo A: listas de verificación individuales}
	
	\begin{figure}[H]
		\centering
		\includegraphics[width=0.85\textwidth]{AnexoA_Diaz.png}
		\caption{Verifacion individual Checklist pag1.}

	\end{figure}
	
	\begin{figure}[H]
		\centering
		\includegraphics[width=0.85\textwidth]{AnexoA_Diaz1.png}
		\caption{Verifacion individual Checklist pag2.}
		
	\end{figure}
	
	\begin{figure}[H]
		\centering
		\includegraphics[width=0.85\textwidth]{AnexoA_Diaz2.png}
		\caption{Verifacion individual Checklist pag3.}
		
	\end{figure}
	
	La demas evidencias de la lista de verificación individual se encuentran el repositorio.
	
	\section{Anexo B: registro de defectos y acta Fagan}

	
	\begin{figure}[H]
		\centering
		\includegraphics[width=1.10\textwidth]{Defectos_Fagan.png}
		\caption{Registro de defectos consolidados.}
		
	\end{figure}
	
	
	
	
	\section{Anexo C: RFC y acta del CCB}

	
	\begin{figure}[H]
		\centering
		\includegraphics[width=1.10\textwidth]{Acta_CCB-firma.png}
		\caption{Acta firmada del Change control board.}
		
	\end{figure}
	
	\begin{figure}[H]
		\centering
		\includegraphics[width=1.10\textwidth]{Acta_CCB (2).png}
		\caption{Acta firmada del Change control board.}
		
	\end{figure}
	

	
	Las instancias de los formularios de la RFC se adjuntan completas en el repositorio. 
	
	\section{Anexo D: exportación CSV}

	\begin{figure}[H]
		
	\centering
		\includegraphics[width=1.10\textwidth]{Exportacion.png}
		\caption{Acta firmada del Change control board.}
	
	\end{figure}
	
	\section{Anexo E: evidencias CASE y Git}
	Insertar: tablero completo, tres issues con enlaces, panel estadístico, \texttt{git log --oneline --graph --decorate} y \texttt{git tag -n}. Todas las capturas deben mostrar proyecto y cuenta del equipo.
	
	\begin{figure}[H]
		\centering
		\includegraphics[width=1.10\textwidth]{Git_Code.png}
		\caption{Acta firmada del Change control board.}
		
	\end{figure}
	
	\section{Anexo F: declaración de uso de inteligencia artificial}
	\begin{table}[H]
		\centering
		\caption{Declaración transparente de asistencia con IA}
		\begin{tabularx}{\textwidth}{p{3.5cm}X}
			\toprule
			Herramienta y versión & OpenAI ChatGPT/Codex, versión disponible el 4 de agosto de 2026.\\
			Tarea asistida & Análisis de estructura de la guía, propuesta de organización \LaTeX, revisión preliminar del repositorio y apoyo de redacción.\\
			Fragmentos afectados & Plantilla del informe, fundamento teórico, tablas comparativas, checklists, RFC, matriz, métricas y exportación estructurada del backlog.\\
			Verificación crítica & El equipo debe contrastar cada afirmación con las normas citadas, el ERS, la matriz, los commits y las evidencias reales. Los candidatos no se convierten en defectos, RFC ni decisiones hasta ser validados y firmados por el equipo.\\
			Limitaciones & La IA no participó en la sesión Fagan ni en el CCB, no produjo capturas auténticas y no puede firmar ni acreditar aportes.\\
			Firmas &  Diaz Pontón Steven - Tigasi Sampedro Paúl - Vera Gomez Anthony\\
			\bottomrule
		\end{tabularx}
	\end{table}
	

\end{document}
